% Copyright (C) 2026 Yu Zhai <me@zhaiyusci.net>
% This work may be distributed and/or modified under the LaTeX Project
% Public License, version 1.3c or (at your option) any later version.
% https://www.latex-project.org/lppl.txt
% LPPL maintenance status: maintained.
% Current Maintainer: Yu Zhai <me@zhaiyusci.net>.
% The distributed Work is listed in MANIFEST.md.
%
\documentclass[a4paper,11pt]{article}
\usepackage[T1]{fontenc}
\usepackage{lmodern}
\usepackage[margin=27mm,headheight=14pt]{geometry}
\usepackage{booktabs,array,xcolor,listings,fancyhdr}
\usepackage{matchedsplit}
\usepackage{amssymb}
\usepackage[colorlinks=true,linkcolor=blue!45!black,urlcolor=blue!45!black]{hyperref}
\hypersetup{pdftitle={The matchedsplit package},pdfauthor={Yu Zhai},pdfsubject={Matched delimiters across manually broken equations}}
\lstset{language=[LaTeX]TeX,basicstyle=\ttfamily\small,
 columns=fullflexible,keepspaces=true,breaklines=true,
 frame=single,rulecolor=\color{black!20},framesep=6pt,
 aboveskip=10pt,belowskip=10pt,showstringspaces=false}
\newcommand\pkg[1]{\textsf{#1}}
\newcommand\env[1]{\texttt{#1}}
\newcommand\cs[1]{\texttt{\textbackslash#1}}
\pagestyle{fancy}
\fancyhf{}
\fancyhead[L]{\small\pkg{matchedsplit}}
\fancyhead[R]{\small User manual, v0.8.0}
\fancyfoot[C]{\thepage}
\setlength\parindent{0pt}
\setlength\parskip{5pt plus 1pt}
\setcounter{tocdepth}{1}
\title{The \pkg{matchedsplit} package\\[5pt]
\large Matched delimiters across manually broken equations}
\author{Yu Zhai\\\small\href{mailto:me@zhaiyusci.net}{\texttt{me@zhaiyusci.net}}}
\date{Version 0.8.0\quad---\quad 29 October 2026}
\begin{document}
\maketitle
\thispagestyle{empty}
\begin{abstract}
\noindent
The \pkg{matchedsplit} package permits a \cs{left}, \cs{middle},
\cs{right} construction to cross manually chosen alignment rows.
Each delimiter pair is sized independently from its contents, including
nested pairs. The package provides an unnumbered inner alignment and
numbered or unnumbered display alignments. Optional package settings
replace the \env{split} and \env{align} environments of \pkg{amsmath}.
\end{abstract}
\tableofcontents
\vfill
\textcopyright\ 2026 Yu Zhai. Distributed under the LaTeX Project Public License,
version 1.3c or later. Maintenance status: maintained. Current maintainer:
Yu Zhai, \href{mailto:me@zhaiyusci.net}{\texttt{me@zhaiyusci.net}}.

\newpage
\section{Getting started}
\subsection{Requirements and installation}
The package requires \LaTeXe{} with the modern document-command,
environment-copy and hook interfaces, and loads \pkg{amsmath}, \pkg{xcolor} and \pkg{l3keys2e} automatically.
Testing used the pdf\LaTeX{}, Xe\LaTeX{} and Lua\LaTeX{} engines supplied
with TeX Live 2026. Earlier formats have not been verified.

For a local installation, place \texttt{matchedsplit.sty} beside the main
source file. For a personal TEXMF installation, place it in
\path{tex/latex/matchedsplit/} below your local TEXMF tree and refresh
the filename database if your installation requires it.
\begin{lstlisting}
\usepackage{matchedsplit}
\end{lstlisting}
No package option is needed to use the new environments. Loading the
package without options leaves the native \env{split}, \env{align} and
\env{align*} definitions unchanged.

\subsection{A complete example}
\begin{lstlisting}
\documentclass{article}
\usepackage{matchedsplit}
\begin{document}
\begin{matchedalign}
F &= \left( a+\frac{x}{y} \label{eq:first}\\
  &\quad +\frac{p}{q}\right)\label{eq:second}
\end{matchedalign}
See \eqref{eq:second}.
\end{document}
\end{lstlisting}
The two rows receive separate numbers, while the parentheses retain the
same size:
\begin{matchedalign}
F &= \left( a+\frac{x}{y} \label{manual:first}\\
  &\quad +\frac{p}{q}\right)\label{manual:second}
\end{matchedalign}
Compile twice to resolve references, such as \eqref{manual:second}.
The package does not choose line breaks: the author supplies \verb|\\|
and alignment markers \verb|&|.

\newpage
\section{Environments and package options}
\subsection{Environment reference}
\begin{center}
\begin{tabular}{@{}ll@{}}
\toprule
Environment & Numbering and placement\\
\midrule
\env{matchedsplit} & Inner math alignment; numbering belongs to its parent.\\
\env{matchedalign} & Display alignment; one number per row by default.\\
\env{matchedalign*} & Display alignment; no automatic numbers.\\
\bottomrule
\end{tabular}
\end{center}
Use \env{matchedsplit} inside \env{equation}, \env{equation*},
\verb|\[...\]|, or an appropriate outer math alignment. Use
\env{matchedalign} directly in document text, not inside \env{equation}
or \verb|\[...\]|. The special \pkg{amsmath} arrangement of an
\env{align} inside \env{gather} has also been tested with the replacement.

\subsection{One number for the complete formula}
\begin{lstlisting}
\begin{equation}\label{eq:whole}
\begin{matchedsplit}
G &= \left(\frac{a}{b}\\
  &\quad +c\right)
\end{matchedsplit}
\end{equation}
\end{lstlisting}
\begin{equation}\label{manual:whole}
\begin{matchedsplit}
G &= \left(\frac{a}{b}\\ &\quad +c\right)
\end{matchedsplit}
\end{equation}

\subsection{Replacing the native environments}
\begin{center}
\begin{tabular}{@{}lll@{}}
\toprule
Option & Default & Effect\\
\midrule
\texttt{overridesplit} & Off & Replace \env{split}.\\
\texttt{overridealign} & Off & Replace \env{align} and \env{align*}.\\
\texttt{rainbow} & Off & Color delimiters by nesting depth.\\
\bottomrule
\end{tabular}
\end{center}
\begin{lstlisting}
\usepackage[overridesplit,overridealign]{matchedsplit}
\end{lstlisting}
The options are independent. The replacement \env{split} uses an
\env{aligned} layout; it does not reproduce every width calculation of
native \env{split}. The replacement \env{align} delegates numbering and
layout to the saved native \pkg{amsmath} implementation.

Configure \pkg{amsmath} before loading this package, for example:
\begin{lstlisting}
\usepackage[tbtags,leqno]{amsmath}
\usepackage[overridesplit]{matchedsplit}
\end{lstlisting}
With \texttt{tbtags}, replacement \env{split} aligns a right-hand tag with
the last row, or a left-hand tag with the first row. Otherwise it centers
the block vertically for tagging.

\newpage
\section{Delimiter sizing and nested structures}
A pair is sized from its mathematical contents after removing the
outer alignment markers, row breaks and supported numbering commands.
The size is \emph{not} the total height of the multirow display. A tall
fraction affects every pair that contains it, but not an unrelated pair.

\subsection{Nested pairs}
\begin{lstlisting}
\begin{matchedalign}
F &= \left( a+\left[\frac{x}{y}+b \\
  &\quad +c\right]+d \\
  &\quad +\frac{p}{q}\right)
\end{matchedalign}
\end{lstlisting}
\begin{matchedalign}
F &= \left( a+\left[\frac{x}{y}+b \\
  &\quad +c\right]+d \\
  &\quad +\frac{p}{q}\right)
\end{matchedalign}
The square brackets form the inner pair and are measured separately
from the surrounding parentheses.

\subsection{Middle and invisible delimiters}
\cs{middle} belongs to the innermost currently open pair. A period is
accepted as an invisible delimiter, as in ordinary \TeX{}.
\begin{lstlisting}
\begin{matchedalign*}
S &= \left\{ x\in\mathbb{R} \middle| \frac{x}{2}>0 \\
  &\quad \middle| x<10\right\}
\end{matchedalign*}
\end{lstlisting}
\begin{matchedalign*}
S &= \left\{ x\in\mathbb{R} \middle| \frac{x}{2}>0 \\
  &\quad \middle| x<10\right\}
\end{matchedalign*}
Use standard single-token delimiter characters or commands, for example
\verb|(|, \verb|\{|, \verb|\langle| and \verb|.|.

\subsection{Groups and inner environments}
Braced groups retain their braces and are not recursively rewritten.
Inner environments, including matrices, retain their own alignment
markers and row breaks. A delimiter pair must not cross a group or
inner-environment boundary. Native math inside such a structure must
already be valid on its own; the outer scanner does not repair it.

\newpage
\section{Row control, tags and references}
Supported row endings are \verb|\\|, \verb|\\*|,
\verb|\\[3pt]| and \verb|\\*[3pt]|. Their spacing arguments are
retained for layout and excluded from delimiter measurements.

In \env{matchedalign}, use \cs{label} and \cs{eqref} as usual.
\cs{notag} or \cs{nonumber} suppresses a row number.
\verb|\tag{A}| supplies a custom tag; \verb|\tag*{B}| omits its
usual surrounding parentheses. The counter behavior is that of
\pkg{amsmath}, including the effect of custom tags on automatic numbering.
\begin{lstlisting}
\begin{matchedalign}
H &= \left(\frac{a}{b}\tag{A}\label{eq:custom}\\
  &\quad +x\notag\\
  &\quad +y\right)\label{eq:last}
\end{matchedalign}
\end{lstlisting}
\begin{matchedalign}
H &= \left(\frac{a}{b}\tag{A}\label{manual:custom}\\
  &\quad +x\notag\\
  &\quad +y\right)\label{manual:last}
\end{matchedalign}
The custom tag can be referenced as \eqref{manual:custom}.
Labels and tags must be at the top level of the relevant row and use
braced arguments. Commands hidden in braced groups are not extracted
from measurement content.

\subsection{Text between rows}
\cs{intertext} is supported in the display alignments, including while
a delimiter pair remains open:
\begin{lstlisting}
\begin{matchedalign}
f &= \left\langle\frac{a}{b}\\
\intertext{The same delimiter pair continues.}
  &\quad +c\right\rangle
\end{matchedalign}
\end{lstlisting}
\begin{matchedalign}
f &= \left\langle\frac{a}{b}\\
\intertext{The same delimiter pair continues.}
  &\quad +c\right\rangle
\end{matchedalign}
The text is not measured as part of the mathematics.
\verb|\displaybreak[0]| and other native break levels are retained for
\pkg{amsmath}. They do not make a boxed \env{matchedsplit} breakable.
Comprehensive multipage pagination tests have not been performed.

\newpage
\section{Rainbow delimiters}
The optional rainbow mode colors delimiters by nesting depth. Matching
left and right delimiters, and any middle delimiters belonging to that
pair, use the same color across rows. Ordinary formula contents and
numbers retain their surrounding color. This feature is inspired by
rainbow-parenthesis highlighting in editors such as Vim's niji.

\subsection{Package defaults and local overrides}
\begin{lstlisting}
\usepackage[rainbow]{matchedsplit}
% Enable globally, but disable for this display:
\begin{matchedalign}[rainbow=false]
f &= \left(a\\ &+b\right)
\end{matchedalign}
\end{lstlisting}
Without a package default, opt in for a single environment:
\begin{lstlisting}
\begin{matchedalign*}[rainbow]
f &= \left\{\frac{a}{b}+\left[x+\left(y\middle|z\\
  &\quad +1\right)\right]\right\}
\end{matchedalign*}
\end{lstlisting}
\begin{matchedalign*}[rainbow]
f &= \left\{\frac{a}{b}+\left[x+\left(y\middle|z\\
  &\quad +1\right)\right]\right\}
\end{matchedalign*}
The optional argument is available on \env{matchedsplit},
\env{matchedalign} and \env{matchedalign*}. With the corresponding
replacement options enabled it is also available on \env{split},
\env{align} and \env{align*}. Environment settings override package
defaults for that environment only; they do not change later displays.
An omitted value in \texttt{rainbow} means \texttt{rainbow=true}.

\subsection{Choosing a palette}
\begin{lstlisting}
\usepackage[rainbow,colors={blue,orange,teal}]{matchedsplit}
% A different palette for one display:
\begin{matchedalign}[rainbow,colors={red!70!black,blue}]
f &= \left(a+\left[b\\ &+c\right]\right)
\end{matchedalign}
\end{lstlisting}
The \texttt{colors} key accepts a nonempty, braced comma-separated list
of \pkg{xcolor} names or expressions. It does not enable rainbow mode
by itself. Undefined colors are diagnosed by \pkg{xcolor}. Configure
\pkg{xcolor} before loading this package if additional color-name sets
are required.

The default palette is blue, orange, green, purple, red and cyan, using
the package's \texttt{msblue}, \texttt{msorange}, \texttt{msgreen},
\texttt{mspurple}, \texttt{msred} and \texttt{mscyan} definitions.
Colors cycle when nesting exceeds the palette length. Sibling pairs at
the same depth reuse the same color, and every environment starts at
depth one. Invisible delimiters still participate in depth counting.
Only directly parsed delimiters are colored: groups, native inner
environments and delimiters hidden in macros remain outside the scanner.

\newpage
\section{Compatibility, diagnostics and limitations}
\subsection{Scope of support}
The implementation targets display-style mathematics with explicit,
manual row breaks. It is not a general inline-math or script-style
formatter. User macros are not expanded to discover hidden delimiters,
alignment markers or metadata. Nested \env{matchedsplit} environments
are outside the supported interface.

Measurement executes mathematical contents more than once, especially
with nested pairs. Avoid counters, file writes or other side effects in
formula contents. Each middle delimiter is rendered as a relation atom;
left and right delimiters are opening and closing atoms. Exact horizontal
spacing need not match an unsplit native construction in every case.

The package uses internal \pkg{amsmath} bindings. Native \env{alignat},
\env{xalignat}, \env{xxalignat} and \env{flalign}, including their defined
starred variants, retain their native behavior and do not acquire
cross-row delimiter support. Their coexistence with the replacement
has been tested. No general compatibility claim is made for third-party
packages such as \pkg{mathtools}, \pkg{hyperref} or \pkg{unicode-math},
or for publisher-specific document classes. This manual uses
\pkg{hyperref}; that is an example configuration, not a complete
compatibility test suite.

\subsection{Diagnostics}
The package reports unmatched \cs{right} or \cs{middle}, unclosed
\cs{left}, a missing final delimiter, missing braced arguments for
\cs{label}, \cs{tag} or \cs{intertext}, and missing closing brackets in
row-spacing arguments. Invalid native mathematics and invalid delimiter
tokens may instead produce a \TeX{} or \pkg{amsmath} error.

Since an environment body is collected before processing, an error may
be reported at its closing \verb|\end| rather than at the originating row.

\subsection{Implementation and validation}
A stack tracks open pairs. Each pair is measured in display style with
invisible boundaries. Local delimiter constructions with height/depth
phantoms allow its ends to occupy different alignment cells. Parser state
and box assignments are local to the environment. Repeated token-list
copying and measurement can be expensive for long or deeply nested input.

The regression suite checks dimensions against native \TeX{}, label
values, option combinations and failure diagnostics. All 45 positive
engine/case combinations and twelve pdf\LaTeX{} negative cases pass in the
tested configuration. Selected PDF outputs have also been inspected.
\newpage
\section{Version history and distribution information}
\begin{description}
\item[0.8.0 (2026-10-29)] Add package and environment rainbow options,
local palette overrides and depth-based delimiter colors.
\item[0.7.1 (2026-10-29)] Fix intertext measurement and align inside gather;
validate metadata arguments and row-spacing delimiters. Unify the public
package name as \pkg{matchedsplit}.
\item[0.7] Add numbered display environments and optional align replacement.
\item[0.6] Add optional split replacement.
\item[0.5] Repair nested-pair measurement and preservation of token groups.
\end{description}
This package is supplied as a directly maintained \texttt{.sty} file;
no docstrip extraction is required. The documentation source is
\texttt{matchedsplit.tex}. Copyright \textcopyright\ 2026 Yu Zhai. This work may be distributed and/or
modified under the LaTeX Project Public License, version 1.3c or (at your
option) any later version. Its maintenance status is \emph{maintained};
the current maintainer is Yu Zhai. See \texttt{LICENSE} and
\texttt{MANIFEST.md} for the license notice and scope of the work.
Send bug reports with a minimal example, engine version and log to
\href{mailto:me@zhaiyusci.net}{\texttt{me@zhaiyusci.net}}.
\end{document}

